\documentclass[../../main.tex]{subfiles}

\graphicspath{{\subfix{images/}}}

\begin{document}

\chapter{Implementarea Modulelor de Inteligență Artificială}
\label{ch:implementare_ai}

\section{Modulul de recunoaștere facială}
\label{sec:modul_facial}

Modulul de recunoaștere facială constituie componenta centrală a sistemului de
control al accesului și este implementat în clasa \inlinecode{FacialRecognitionSystem}
din fișierul \inlinecode{facial\_recognition\_system.py}. Rolul său este de a
identifica, în fiecare cadru video, persoanele înregistrate anterior în baza de
date și de a distinge între persoanele cunoscute (angajați) și cele necunoscute.
Spre deosebire de pipeline-urile clasice de recunoaștere facială, care înlănțuie
un detector de fețe, un model separat de extragere a embeddings-urilor și încă
unul pentru estimarea atributelor demografice, implementarea de față consolidează
toate aceste sarcini într-o singură trecere prin pachetul de modele
\textbfit{buffalo\_s} al framework-ului InsightFace, ale cărui fundamente
teoretice au fost prezentate în secțiunea~\ref{sec:recunoastere_faciala}.

Pipeline-ul intern al modulului are patru etape, descrise în subcapitolele care
urmează: (1)~detecția fețelor și extragerea embeddings-urilor cu InsightFace,
(2)~indexarea și căutarea vectorilor cu FAISS, (3)~strategia de stabilire a
identității pe baza distanței și (4)~optimizarea prin omiterea analizei
demografice pentru fețele deja recunoscute. La nivelul aplicației, modulul rulează
într-un fir de execuție dedicat (\textit{Thread~2}), care preia cadre dintr-o coadă
de tip \inlinecode{queue.Queue}, apelează metoda \inlinecode{process\_frame()} și
publică rezultatele în coada de fuziune (a se vedea secțiunea~\ref{sec:fuziune}).

\subsection{Detecția fețelor și extragerea embeddings-urilor}
\label{subsec:detectie_extragere}

Detecția fețelor și extragerea descriptorilor numerici (embeddings) sunt realizate
de obiectul \inlinecode{FaceAnalysis} din InsightFace, inițializat cu pachetul de
modele \inlinecode{buffalo\_s}. Acest pachet reunește trei submodele care, în
abordarea tradițională, ar fi necesitat încărcarea și apelarea separată:

\begin{itemize}
    \item \textbf{det\_500m} -- detectorul de fețe, care returnează bounding box-ul
    și cele cinci puncte caracteristice (landmark-uri) pentru fiecare față din cadru;
    \item \textbf{w600k\_mbf} -- rețeaua de recunoaștere, o variantă MobileFaceNet
    antrenată pe datasetul Glint360K cu funcția de pierdere ArcFace, care produce
    embedding-ul facial de 512 dimensiuni;
    \item capul \textbf{genderage} -- care estimează vârsta și genul persoanei.
\end{itemize}

Alegerea variantei \inlinecode{buffalo\_s} (în locul, de exemplu, a variantei
\inlinecode{buffalo\_sc}) este motivată tehnic: rețeaua de recunoaștere
\inlinecode{w600k\_mbf} este identică în ambele pachete, astfel încât embedding-urile
de 512 dimensiuni rămân compatibile, însă \inlinecode{buffalo\_s}  include modulul suplimentar \texttt{genderage}, indispensabil pentru analiza demografică
(secțiunea~\ref{sec:analiza_demografica}). În lipsa acestui modul de predicție, estimările privind vârsta și genul ar fi indisponibile, sistemul returnând valoarea \texttt{None}.

Pentru fiecare cadru, metoda \inlinecode{process\_frame()} apelează
\inlinecode{self.face\_app.} \inlinecode{get(frame)}, care efectuează într-o singură trecere
detecția tuturor fețelor și, pentru fiecare, calculează bounding box-ul,
landmark-urile, embedding-ul L2-normalizat (\inlinecode{normed\_embedding}, vector
de 512 de elemente de tip \inlinecode{float32}) și estimările de vârstă și gen.
Bounding box-ul de tip \inlinecode{[x1, y1, x2, y2]} returnat de InsightFace este
convertit în formatul \inlinecode{(x, y, w, h)} și ancorat în limitele cadrului de
metoda statică \inlinecode{\_bbox\_xyxy\_to\_xywh()}, pentru a evita decupaje în
afara imaginii.

\subsubsection{Configurarea execuției pe GPU}
Inferența InsightFace rulează pe GPU prin runtime-ul ONNX, providerul \linebreak
\inlinecode{CUDAExecutionProvider} fiind configurat explicit cu opțiunea
\inlinecode{cudnn\_conv\_algo\_} \inlinecode{search="HEURISTIC"}. Această opțiune este o decizie de implementare importantă în contextul arhitecturii multi-threaded a aplicației:
strategia implicită, \inlinecode{EXHAUSTIVE}, măsoară empiric timpul fiecărui
algoritm de convoluție folosind mecanismul de \textit{stream capture} al CUDA. Cum
InsightFace rulează concurent cu modelele PyTorch (YOLO pentru plăcuțe, foc și arme,
respectiv SlowFast pentru acțiuni) pe același GPU, o astfel de captură în desfășurare
ar întrerupe orice kernel PyTorch paralel cu eroarea \textit{``operation not permitted
when stream is capturing''}. Varianta \inlinecode{HEURISTIC} selectează algoritmul pe
baza unor euristici, fără benchmarking și deci fără captură, eliminând coliziunile
între firele de execuție.

\subsubsection{Gestionarea bibliotecilor CUDA și a TensorFlow}
La importul modulului, funcția \inlinecode{\_preload\_nvidia\_pypi\_libs()}
preîncarcă explicit, prin \inlinecode{ctypes.CDLL(..., RTLD\_GLOBAL)}, bibliotecile
partajate CUDA, cuDNN și cuBLAS livrate în pachetele pip \inlinecode{nvidia-*-cu12},
astfel încât \linebreak \inlinecode{onnxruntime} să le poată găsi fără a necesita setarea
manuală a variabilei \linebreak \inlinecode{LD\_LIBRARY\_PATH}. În plus, TensorFlow (folosit
de modelul de emoții) este forțat să nu folosească GPU-ul prin
\inlinecode{set\_visible\_devices([], "GPU")}, deoarece menținerea lui pe GPU
provoacă conflicte de context CUDA cu \inlinecode{onnxruntime} \inlinecode{-gpu} și segmentation
fault la inițializare.

\subsubsection{Înregistrarea persoanelor}
Embedding-urile persoanelor cunoscute sunt generate în momentul înregistrării acestora în sistem. Metoda \inlinecode{register\_person()} primește o listă de imagini faciale și execută modelul InsightFace pe fiecare dintre ele. În situația în care sunt detectate mai multe fețe într-o imagine, sistemul o selectează pe cea cu aria maximă a casetei de încadrare (\emph{bounding box}), prin intermediul funcției \inlinecode{\_extract\_single\_face()}. Ulterior, este extras \emph{embedding}-ul de 512 dimensiuni, care va fi stocat atât în baza de date, cât și în indexul FAISS. Pentru încărcările de imagini efectuate prin interfața web, metoda \inlinecode{extract\_embedding\_} \inlinecode{from\_image()} acceptă exclusiv imaginile care conțin o singură față detectabilă, prevenind astfel riscul asocierii ambigue a unui \emph{embedding} cu o persoană.

\subsection{Indexarea cu FAISS FlatL2}
\label{subsec:indexare_faiss}

Căutarea identității unei fețe presupune compararea embedding-ului ei cu toți
descriptorii persoanelor înregistrate. Pentru a face această operație eficientă,
sistemul folosește biblioteca FAISS, încapsulată în clasa \inlinecode{FaissIndex}
din fișierul \inlinecode{faiss\_index.py}. Tipul de index utilizat este
\textbfit{IndexFlatL2}, care realizează o căutare \textit{exactă} (\textit{brute-force})
folosind distanța euclidiană L2 între vectorul de interogare și fiecare vector din
index. Spre deosebire de indecșii aproximativi (precum IVFFlat), căutarea exactă nu
necesită o etapă de antrenare și garantează găsirea celui mai apropiat vecin; aceasta
este alegerea potrivită pentru dimensiunea tipică a unei baze de angajați, unde numărul
de embedding-uri este suficient de mic încât căutarea liniară să rămână instantanee.

\subsubsection{Maparea identităților}
FAISS operează cu indici poziționali interni (0, 1, 2, \dots), nu cu identificatorii
persoanelor din baza de date. De aceea, clasa \inlinecode{FaissIndex} menține două
dicționare de mapare: \inlinecode{id\_map}, care asociază fiecărui index FAISS un
\inlinecode{person\_id}, și \inlinecode{reverse\_id\_map}, care leagă fiecare
\inlinecode{person\_id} de lista de indici corespunzători (o persoană poate avea mai
multe embedding-uri, câte unul pentru fiecare imagine înregistrată).

\subsubsection{Persistența și reconstrucția indexului}
Indexul FAISS este o structură de date alocată în memorie, în timp ce baza de date PostgreSQL constituie punctul central de stocare și referință pentru toți vectorii de \emph{embedding}. În tabelul \inlinecode{face\_embeddings}, fiecare \emph{embedding} este salvat în coloana \inlinecode{embedding} de tip \inlinecode{BYTEA}, fiind serializat cu ajutorul bibliotecii \inlinecode{pickle}. La pornirea sistemului, metoda \inlinecode{\_load\_embeddings\_from\_db()} preia toate datele prin \inlinecode{get\_all\_embeddings()} și reconstruiește complet indexul prin apelarea funcției \inlinecode{rebuild\_index()}. Această procedură de reconstrucție este declanșată automat ori de câte ori o persoană este înregistrată, modificată sau ștearsă, garantând astfel consistența strictă între stocarea persistentă și indexul volatil. 

\subsubsection{Siguranța accesului concurent}
Întrucât firul de recunoaștere facială interoghează indexul în fiecare cadru, în
timp ce firele care deservesc cererile HTTP îl pot reconstrui (la înregistrarea sau
ștergerea persoanelor), iar operațiile \inlinecode{add}/\inlinecode{search} ale FAISS
nu sunt sigure pentru execuția concurentă, toate accesările indexului și ale
dicționarelor de mapare sunt serializate cu un lock reentrant
(\inlinecode{threading.RLock}). Caracterul reentrant este necesar deoarece
\inlinecode{rebuild\_index()} apelează intern \inlinecode{add\_embeddings} \inlinecode{\_batch()}, care preia același lock.

\subsection{Strategia de matching}
\label{subsec:matching}

Stabilirea identității unei fețe pe baza distanței euclidiene este implementată în
metoda \inlinecode{\_match\_embedding()}. Aceasta interoghează mai întâi indexul
pentru cel mai apropiat vecin (\inlinecode{k=1}) cu un prag ``infinit'', astfel încât
distanța reală până la cel mai apropiat vector să poată fi înregistrată pentru
calibrare (activabilă prin variabila de mediu \inlinecode{FR\_DEBUG}), după care aplică
manual pragul de decizie configurat.

\subsubsection{De la distanța L2 la similaritatea cosinus}
Deoarece embeddings-urile sunt L2-normalizate, între distanța euclidiană la pătrat și
similaritatea cosinus a doi vectori există relația exactă
$L2^2 = 2 - 2\cos\theta$, echivalentă cu $\cos\theta = 1 - L2^2/2$. Această echivalență
permite raportarea identității prin două mărimi interpretabile:

\begin{itemize}
    \item \inlinecode{recognition\_threshold = 1.0} -- pragul aplicat distanței L2.
    Un prag $L2 = 1{,}0$ corespunde unei similarități cosinus de $0{,}5$. Empiric,
    distribuția distanțelor pentru aceeași persoană se situează în intervalul
    $[0{,}6;\ 0{,}85]$ pe camera web folosită, deci un candidat cu $L2 \geq 1{,}0$
    este respins ca necunoscut;
    \item \inlinecode{min\_confidence = 0.4} -- pragul aplicat încrederii, exprimată
    ca similaritate cosinus. Valoarea $0{,}4$ este punctul de operare clasic
    ``aceeași persoană'' pentru ArcFace.
\end{itemize}

Dacă distanța celui mai apropiat vecin depășește \inlinecode{recognition\_threshold},
metoda returnează \inlinecode{None} (față necunoscută). În caz contrar, distanța este
convertită în similaritate cosinus, mărginită la intervalul $[0;\ 1]$ și raportată ca
nivel de încredere; identificatorul persoanei este apoi rezolvat în datele complete
(nume, ID angajat, departament) prin \inlinecode{get\_person\_by\_id()}. În metoda
\inlinecode{process\_frame()}, o față este considerată cunoscută doar dacă există o
potrivire \textit{și} încrederea ei depășește \inlinecode{min\_confidence}. Pentru
fiecare potrivire validă se înregistrează un eveniment de acces în baza de date prin
\inlinecode{log\_access()}.

Vizual, fețele cunoscute sunt marcate cu un chenar verde și eticheta
\inlinecode{nume} \inlinecode{(încredere)}, iar cele necunoscute cu un chenar portocaliu și eticheta
\inlinecode{Unknown}, prin metoda \inlinecode{\_draw\_face()}.

\subsection{Optimizarea: omiterea analizei demografice pentru fețele cunoscute}
\label{subsec:optimizare_demografica}

O optimizare deliberată la nivelul \emph{pipeline}-ului de procesare facială constă în condiționarea analizei demografice de rezultatul etapei de recunoaștere. În cadrul funcției \inlinecode{process\_frame()}, atributele demografice (precum vârsta și genul, generate de modulul \inlinecode{genderage} al arhitecturii InsightFace, alături de emoția prezisă de rețeaua Mini-Xception, detaliată în Secțiunea~\ref{sec:analiza_demografica}) sunt calculate \emph{exclusiv pentru fețele necunoscute}. În cazul unei persoane identificate cu succes, câmpurile \inlinecode{age}, \inlinecode{gender}, \inlinecode{emotion} și \inlinecode{emotion\_confidence} sunt inițializate direct cu valoarea \inlinecode{None}.

Raționamentul are atât o componentă de proiectare, cât și una de performanță. Din perspectivă semantică, calcularea informațiilor demografice pentru o persoană deja recunoscută reprezintă o operațiune redundantă, întrucât identitatea, vârsta exactă și departamentul sunt deja disponibile în baza de date. Din punct de vedere al performanței, această decizie elimină necesitatea unei treceri suplimentare (\emph{forward-pass}) prin rețeaua de clasificare a emoțiilor pentru fețele cunoscute, la fiecare cadru. Economia de resurse este considerabilă, în special pentru că inferența rețelei Mini-Xception se execută pe unitatea CPU. Prin urmare, costul computațional aferent analizei demografice este asumat doar atunci când generează informație utilă, adică în cazul persoanelor neidentificate,  pentru care caracterizarea demografică și emoțională devine relevantă în scenariile de supraveghere.

\section{Analiza demografică și emoțională}
\label{sec:analiza_demografica}

Pe lângă identificarea persoanelor cunoscute, sistemul caracterizează demografic și
emoțional persoanele \textit{necunoscute}, adică acele fețe care nu pot fi asociate
niciunei intrări din baza de date. Această analiză produce trei atribute pentru
fiecare astfel de față: vârsta estimată, genul și emoția. Implementarea este integrată
în aceeași clasă \inlinecode{FacialRecognitionSystem} și în aceeași metodă
\inlinecode{process\_frame()} ca recunoașterea facială, evitând o trecere separată de
detecție a fețelor. După cum a fost detaliat în secțiunea~\ref{subsec:optimizare_demografica}, aplicarea analizei demografice este strict condiționată (vizând exclusiv fețele necunoscute), întrucât estimarea acestor atribute devine redundantă în cazul persoanelor deja identificate.

Extragerea atributelor se bazează pe două modele distincte: vârsta și genul sunt furnizate de modulul \inlinecode{genderage} din cadrul arhitecturii InsightFace (deja evaluat implicit prin execuția pachetului \inlinecode{buffalo\_s}), în timp ce starea emoțională este estimată de o rețea convoluțională independentă, Mini-Xception, aplicată direct pe regiunea facială decupată (\emph{face crop}).

\subsection{Estimarea vârstei și genului}
\label{subsec:varsta_gen}

Vârsta și genul nu necesită un model dedicat: ambele sunt generate de modulul
\inlinecode{genderage} al pachetului \inlinecode{buffalo\_s}, în aceeași trecere
InsightFace care furnizează bounding box-ul, landmark-urile și embedding-ul. Pentru
fiecare obiect \inlinecode{Face} returnat de detector, atributele sunt extrase astfel:

\begin{itemize}
    \item \textbf{Vârsta} -- câmpul \inlinecode{face.age} (valoare reală) este rotunjit
    la cel mai apropiat întreg. Dacă atributul lipsește, vârsta este raportată ca
    \inlinecode{None};
    \item \textbf{Genul} -- câmpul \inlinecode{face.gender} este o valoare întreagă în
    convenția InsightFace (\inlinecode{1}~=~masculin, \inlinecode{0}~=~feminin), care este
    tradusă în etichetele textuale \inlinecode{"Man"}, respectiv \inlinecode{"Woman"}.
\end{itemize}

Această abordare nu adaugă cost de inferență suplimentar: estimările sunt deja calculate
de InsightFace, fiind doar citite și formatate.
(secțiunea~\ref{subsec:detectie_extragere}).

\subsection{Recunoașterea emoțiilor cu Mini-Xception}
\label{subsec:emotii_xception}

Spre deosebire de vârstă și gen, emoția este clasificată de un model separat:
Mini-Xception, o rețea neuronală convoluțională compactă furnizată de pachetul Python
\inlinecode{fer} sub forma fișierului de ponderi \inlinecode{emotion\_model.hdf5}. Modelul este
antrenat pe datasetul FER-2013 și clasifică expresia facială într-una din șapte categorii,
ordonate conform ieșirilor softmax ale modelului în tuplul \inlinecode{\_EMOTION\_LABELS}:
\inlinecode{angry}, \inlinecode{disgust}, \inlinecode{fear}, \inlinecode{happy},
\inlinecode{sad}, \inlinecode{surprise} și \inlinecode{neutral}.

Un aspect important al implementării este faptul că API-ul implicit al pachetului
\inlinecode{fer} este \textit{ocolit complet}. Pachetul \inlinecode{fer} efectuează în mod
normal propria sa detecție de fețe cu un clasificator Haar cascade înainte de a clasifica
emoția -- o trecere redundantă, întrucât InsightFace a detectat deja toate fețele din cadru.
În locul acestui flux, modelul Mini-Xception este încărcat direct prin
\inlinecode{tf.keras.models.load\_model()} (cu \inlinecode{compile=False}) și este alimentat
cu crop-ul facial deja delimitat de bounding box-ul InsightFace. Astfel se elimină o
detecție de fețe duplicată și se garantează că emoția este clasificată exact pe regiunea
identificată de pipeline-ul principal.

Predicția propriu-zisă, implementată în metoda \inlinecode{\_classify\_emotion()}, constă într-o etapă de propagare directă (\emph{forward pass}) prin rețea, urmată de aplicarea operației \inlinecode{argmax} pe cele șapte probabilități furnizate de funcția \emph{softmax}. Rezultatul returnat este un tuplu format din eticheta emoției și scorul de confidență asociat. Dimensiunea de intrare necesară modelului este extrasă automat din proprietatea \inlinecode{input\_shape}, asigurând astfel adaptarea dinamică a etapei de preprocesare (detaliată în continuare) la arhitectura specifică rețelei.

% [DIAGRAMĂ SUGERATĂ] Aici s-ar potrivi o diagramă a fluxului analizei pe o față
% necunoscută: cadru → InsightFace (bbox + landmarks + embedding + vârstă/gen) →
% [dacă față necunoscută] crop bbox → preprocesare → Mini-Xception → emoție.
% Ar evidenția cele două modele (genderage vs Mini-Xception) și ramificarea
% condiționată cunoscut/necunoscut.

\subsection{Preprocesarea crop-ului facial}
\label{subsec:preprocesare_crop}

Înainte de a fi evaluat de rețeaua de clasificare a emoțiilor, decupajul facial (extras din imaginea de ansamblu pe baza casetei de încadrare furnizate de InsightFace, în format BGR) parcurge o etapă riguroasă de preprocesare. Acest proces reproduce fidel fluxul intern de prelucrare al pachetului \inlinecode{fer}, garantând astfel faptul că ponderile antrenate ale modelului primesc datele de intrare exact în distribuția statistică necesară. Pașii de transformare, implementați în cadrul metodei \inlinecode{\_classify\_emotion()}, sunt următorii:

\begin{enumerate}
	\item \textbf{Conversia în tonuri de gri}: regiunea facială în format BGR este convertită în scară de gri utilizând funcția \inlinecode{cv2.cvtColor}, deoarece modelul Mini-Xception procesează exclusiv imagini cu un singur canal de culoare;
	\item \textbf{Redimensionarea}: imaginea rezultată este adusă la dimensiunile de intrare impuse de model (la valoarea \inlinecode{\_emotion\_target\_size}, dedusă dinamic din atributul \inlinecode{input\_shape});
	\item \textbf{Normalizarea}: intensitatea pixelilor este scalată din intervalul $[0, 255]$ în intervalul $[0, 1]$ (prin împărțire la $255$), fiind ulterior mapată în intervalul $[-1, 1]$ aplicând transformarea matematică $(x - 0{,}5) \cdot 2$;
	\item \textbf{Extinderea dimensionalității tensorului}: matricea bidimensională este remodelată la forma \inlinecode{(1, H, W, 1)} (adăugându-se dimensiunea lotului de prelucrare, \emph{batch size}, și axa canalului), acesta fiind formatul structural așteptat de modelele Keras.
\end{enumerate}

Întregul lanț este protejat de un bloc \inlinecode{try/except}: dacă, din orice motiv,
crop-ul este gol sau preprocesarea eșuează, metoda returnează \inlinecode{(None, None)}, iar
fața este raportată fără emoție, fără a întrerupe procesarea celorlalte fețe din cadru.

\subsection{Inferența pe CPU}
\label{subsec:inferenta_cpu}

Spre deosebire de InsightFace, care rulează pe GPU, rețeaua Mini-Xception este rulată
deliberat pe CPU. Această decizie este aplicată chiar la importul modulului, unde
TensorFlow este forțat să nu vadă niciun dispozitiv GPU prin
\inlinecode{tf.config.set\_visible\_devices([], "GPU")}.

Motivația este dublă. În primul rând, menținerea TensorFlow pe GPU în paralel cu
\inlinecode{onnxruntime-gpu} (folosit de InsightFace) provoacă conflicte de context CUDA și
segmentation fault la inițializare; izolarea TensorFlow pe CPU elimină complet această
sursă de instabilitate. În al doilea rând, Mini-Xception este o rețea suficient de mică
încât latența inferenței pe CPU să fie neglijabilă, iar emoția se calculează doar pentru
fețele necunoscute (secțiunea~\ref{subsec:optimizare_demografica}); prin urmare, costul pe
CPU este minor, iar GPU-ul rămâne dedicat integral modelelor grele.

\section{Modulul de recunoaștere a plăcuțelor de înmatriculare}
\label{sec:modul_placute}

Modulul de recunoaștere a plăcuțelor de înmatriculare (ANPR -- \textit{Automatic Number
Plate Recognition}) identifică și citește numerele de înmatriculare ale vehiculelor din
fluxul video, le validează formatul, le caută în baza de date și decide dacă vehiculul
este autorizat. Modulul este implementat în clasa \inlinecode{LicensePlateRecognitionSystem}
din fișierul \inlinecode{license\_plate\_recognition\_} \linebreak \inlinecode{system.py} și rulează într-un fir de
execuție dedicat (\textit{Thread~3}).

\emph{Pipeline}-ul combină un detector de obiecte (YOLOv8) cu un motor de recu-noaștere optică a caracterelor (EasyOCR), între care se interpun mai mulți pași de preprocesare a imaginii și, ulterior, o etapă de stabilizare temporală a rezultatului. O provocare majoră, care a influențat semnificativ implementarea, este faptul că plăcuțele apar în cadru la rezoluții reduse (tipic 116--151~px lățime), insuficiente pentru o recunoaștere OCR fiabilă. Din acest motiv, \emph{pipeline}-ul include un pas explicit de mărire a imaginii (\emph{up-scaling}) și o etapă avansată de filtrare și accentuare a detaliilor. Cele opt subcapitole care urmează descriu, în ordinea execuției, etapele acestui flux de prelucrare.

% [DIAGRAMĂ SUGERATĂ] O diagramă de tip pipeline ar ilustra bine întregul flux:
% cadru → YOLOv8 (detecție + bbox) → expandare bbox (padding 0.20) → up-scaling
% (≥250 px) → preprocesare (grayscale/bilateral/CLAHE/unsharp) → EasyOCR →
% [fallback pe crop color dacă conf<0.4] → validare format RO → PlateTracker
% (IoU + votare temporală) → căutare în BD + decizie autorizare.

\subsection{Detecția plăcuțelor cu YOLOv8}
\label{subsec:detectie_placute}

Localizarea plăcuțelor în cadru este realizată cu ajutorul unui model YOLOv8, instanțiat prin intermediul clasei \inlinecode{YOLO} din pachetul \inlinecode{ultralytics}, pe baza fișierului de ponderi \inlinecode{models/license\_plate/best.pt}. Acesta reprezintă o variantă pre-antrenată, dedicată exclusiv detecției plăcuțelor de înmatriculare (având definită o singură clasă, \inlinecode{license\_plate}). Modelul a fost integrat în sistem în forma sa originală, fără a necesita o etapă suplimentară de antrenare (\emph{fine-tuning}) în cadrul acestui proiect. La momentul inițializării, pentru a optimiza viteza de inferență, sistemul verifică disponibilitatea mediului CUDA și, în caz afirmativ, transferă execuția modelului pe unitatea GPU.

Pentru fiecare cadru procesat, metoda \inlinecode{process\_frame()} apelează funcția \inlinecode{predict()} utilizând parametrul \inlinecode{imgsz=640} și un prag de confidență \inlinecode{conf=0.25}, extrăgând astfel casetele de încadrare candidate (\emph{bounding boxes}). Fiecare casetă, definită prin coordonatele \inlinecode{[x1, y1, x2, y2]}, este ulterior \textit{extinsă} cu 20\% în toate direcțiile prin intermediul funcției \inlinecode{expand\_bbox()} (utilizând un coeficient \inlinecode{pad\_ratio} de \inlinecode{0.20}). Această margine de siguranță mai generoasă garantează includerea limitelor fizice ale plăcuței alături de un context vizual minim, prevenind astfel decuparea accidentală a caracterelor aflate la extremități înaintea evaluării OCR. Regiunea decupată rezultată este transmisă etapelor ulterioare de prelucrare.

\subsection{Up-scaling-ul plăcuțelor mici}
\label{subsec:upscaling}

Deoarece plăcuțele detectate prezintă frecvent dimensiuni insuficiente pentru o recunoaștere OCR precisă, regiunea decupată este procesată de funcția \linebreak \inlinecode{upscale\_if\_small()}, având ca scop redimensionarea la o lățime țintă de \inlinecode{250}~px. Algoritmul aplică următoarea regulă: dacă lățimea curentă se situează sub acest prag, se calculează un factor de scalare $\mathit{scale} = 250 / w$. Acesta este limitat superior la valoarea \inlinecode{3.0} pentru a preveni amplificarea excesivă a zgomotului de imagine. Redimensionarea propriu-zisă utilizează interpolarea bicubică (\inlinecode{cv2.INTER\_CUBIC}), selectată datorită calității superioare a reconstrucției detaliilor în procesul de suprascalare. Casetele care îndeplinesc deja criteriul de lățime minimă rămân nemodificate. Printr-o astfel de operațiune, o regiune inițială de $116 \times 33$~px este transformată într-o imagine de aproximativ $250 \times 70$~px, rezoluție la care performanța algoritmului OCR atinge un nivel optim de fiabilitate.

\subsection{Pipeline-ul de preprocesare pentru OCR}
\label{subsec:preprocesare_ocr}

Imaginea mărită este ulterior pregătită pentru etapa de recunoaștere optică a textului (OCR) prin intermediul funcției \inlinecode{preprocess\_plate()}. Aceasta aplică o secvență de patru operații, optimizată special pentru a maximiza performanța algoritmului EasyOCR în identificarea plăcuțelor de înmatriculare:

\begin{enumerate}
    \item \textbf{Conversie în tonuri de gri} -- imaginea color este transformată în
    grayscale;
    \item \textbf{Filtru bilateral} (\inlinecode{cv2.bilateralFilter} cu $d=11$ și sigma de
    culoare și de spațiu de $17$) -- reduce zgomotul, păstrând în același timp muchiile
    caracterelor;
    \item \textbf{Egalizare locală de histogramă CLAHE} (\inlinecode{clipLimit=2.0},
    \inlinecode{tileGrid} \inlinecode{Size=(8, 8)}) -- compensează umbrele și iluminarea neuniformă;
    \item \textbf{Unsharp mask} -- accentuează muchiile caracterelor după mărire, construit
    ca diferență ponderată între imaginea originală (coeficient $1{,}5$) și o versiune blurată
    Gaussian cu $\sigma = 1{,}0$ (coeficient $-0{,}5$).
\end{enumerate}

O decizie deliberată este aceea de \textit{a nu binariza} imaginea: modelul CRNN folosit de
EasyOCR oferă rezultate mai bune pe imagini continue în tonuri de gri decât pe imagini cu prag
aplicat. Pipeline-ul se oprește, prin urmare, la o imagine grayscale curățată și accentuată,
nu la o imagine alb-negru.

\subsection{Apelul EasyOCR}
\label{subsec:apel_easyocr}

Citirea propriu-zisă a caracterelor este realizată de EasyOCR în metoda
\inlinecode{\_run\_ocr()}. Apelul \inlinecode{reader.readtext()} este parametrizat specific
pentru plăcuțe:

\begin{itemize}
    \item \inlinecode{allowlist} restrânge alfabetul recunoscut la literele mari
    \inlinecode{A--Z} și cifrele \inlinecode{0--9}, eliminând caracterele imposibile pe o
    plăcuță;
    \item \inlinecode{text\_threshold=0.6}, \inlinecode{low\_text=0.3} și
    \inlinecode{mag\_ratio=1.5} reglează sensibilitatea detecției de text și mărirea internă a
    EasyOCR;
    \item \inlinecode{paragraph=False} dezactivează gruparea în paragrafe, nepotrivită pentru
    un singur rând de caractere.
\end{itemize}

Întrucât EasyOCR poate returna mai multe fragmente de text (de exemplu, grupul de litere și
cel de cifre separat), fragmentele sunt sortate de la stânga la dreapta după coordonata $x$ a
colțului stânga-sus, apoi sunt curățate (păstrând doar caractere alfanumerice, transformate în
majuscule) și concatenate într-un singur șir. Încrederea raportată este media aritmetică a
încrederilor fragmentelor componente.

\subsection{Fallback pe imaginea color}
\label{subsec:fallback_color}

Preprocesarea agresivă (grayscale, filtrare, CLAHE, unsharp) îmbunătățește citirea în
majoritatea cazurilor, dar în anumite condiții de iluminare poate degrada rezultatul. Pentru
a acoperi aceste situații, pipeline-ul include un mecanism de \textit{fallback}: dacă citirea
de pe imaginea preprocesată este goală sau are o încredere sub \inlinecode{0.4}, OCR-ul este
reluat pe decupajul color original (nepreprocesat). Dacă această a doua citire are o încredere
mai mare, ea o înlocuiește pe prima. Astfel, sistemul nu rămâne blocat pe o preprocesare care,
ocazional, dăunează în loc să ajute.

\subsection{Tracking-ul plăcuțelor între cadre}
\label{subsec:tracking_placute}

Extragerea textului dintr-un singur cadru este o abordare predispusă la erori, deoarece algoritmul OCR poate confunda cu ușurință caractere similare din punct de vedere vizual (precum \inlinecode{O} și \inlinecode{0}, sau \inlinecode{B} și \inlinecode{8}) de-a lungul cadrelor succesive. Pentru a garanta stabilitatea rezultatului, fiecare detecție este monitorizată temporal prin intermediul clasei \inlinecode{PlateTracker}. Asocierea unei noi detecții cu o traiectorie existentă (\emph{track}) se realizează evaluând suprapunerea spațială: se calculează metrica \emph{IoU} (\emph{Intersection over Union}, prin apelul funcției \inlinecode{\_iou()}) între caseta de încadrare curentă și casetele traiectoriilor active. Detecția este atribuită traiectoriei care prezintă valoarea maximă a metricei IoU, cu condiția ca aceasta să depășească pragul minim definit de \inlinecode{iou\_threshold=0.3}. În situația în care nu se identifică nicio corespondență validă, sistemul inițializează o nouă traiectorie, alocându-i un identificator unic.

Fiecare traiectorie menține un istoric al citirilor recente, stocat într-o structură de date liniară cu capacitate limitată (un obiect \inlinecode{deque} configurat cu parametrul \inlinecode{history=10}). Suplimentar, entitatea folosește parametrul de stare \inlinecode{age}, care contorizează numărul de cadre consecutive în care respectiva traiectorie nu a primit actualizări. La procesarea fiecărui cadru, metoda \inlinecode{age\_tracks()} incrementează acest contor de inactivitate pentru toate traiectoriile monitorizate, eliminându-le automat pe cele care depășesc pragul \inlinecode{max\_age=15} cadre. Acest mecanism asigură oprirea urmăririi pentru vehiculele care au părăsit câmpul vizual al camerei.

\subsection{Validarea formatului}
\label{subsec:validare_format}

Înainte de a fi inclusă în procesul decizional de agregare, fiecare citire este validată în raport cu formatul standard al plăcuțelor de înmatriculare din România, prin intermediul expresiei regulate \inlinecode{ROMANIAN\_PLATE\_RE}. Această regulă acceptă configurații formate din una sau două litere pentru indicativul județului, urmate de două sau trei cifre și un grup final de trei litere (de exemplu, \inlinecode{CJ12ABC} sau, specific pentru București, \inlinecode{B123ABC}). În cadrul metodei \inlinecode{PlateTracker.update()}, un rezultat OCR este adăugat în istoricul instanței de urmărire exclusiv dacă îndeplinește cumulativ trei condiții: este un șir nevid, prezintă un scor de confidență de minimum \inlinecode{0.4} \textit{și} respectă strict formatul impus. Rezultatele incorect formatate (provenite, în general, din erori ale modulului OCR) sunt astfel filtrate anterior etapei de evaluare, contribuind la o creștere semnificativă a robusteței rezultatului final.

\subsection{Votarea temporală}
\label{subsec:votare_temporala}

Determinarea rezultatului final pentru o traiectorie se realizează printr-un proces de \textit{selecție majoritară} aplicat pe istoricul citirilor valide. Această logică de agregare, integrată în metoda \inlinecode{PlateTracker.update()}, se declanșează exclusiv după ce instanța de urmărire a înregistrat un minim de \inlinecode{min\_votes=3} citiri. Datele candidate sunt centralizate utilizând o structură de tip \inlinecode{Counter}, în cadrul căreia fiecare text cumulează suma scorurilor de confidență aferente. Șirul de caractere cu scorul total maxim este desemnat câștigător; dacă numărul recunoașterilor individuale care îl susțin atinge pragul \inlinecode{min\_votes}, acesta devine rezultatul consolidat (starea \inlinecode{emitted}) al traiectoriei, având un scor de confidență final calculat ca media aritmetică a scorurilor individuale.

Interogarea bazei de date este inițiată strict după ce o traiectorie generează o citire stabilizată: funcția \inlinecode{lookup\_owner\_by\_plate()} verifică existența unui proprietar, iar validarea unei înregistrări atestă statutul de vehicul autorizat. Plăcuțele aflate în stadiu de evaluare sau care nu figurează în sistem sunt etichetate drept \inlinecode{UNKNOWN} (sau \inlinecode{"Unknown car"}). La nivelul interfeței vizuale, vehiculele autorizate sunt delimitate printr-o casetă verde, în timp ce vehiculele neautorizate sau necunoscute sunt marcate cu roșu. Acest mecanism de agregare temporală garantează că o plăcuță este validată și supusă verificării doar atunci când există un consens pe mai multe cadre consecutive asupra aceluiași text, eliminând astfel cu succes erorile de citire tranzitorii.
\section{Modulul de detecție foc și fum}
\label{sec:modul_foc_fum}

Modulul de detecție a focului și a fumului identifică în fluxul video apariția flăcărilor și a emisiilor de fum, facilitând declanșarea unor alerte timpurii în scenariile de incendiu. Acesta este implementat în clasa \inlinecode{FireDetectionSystem} din fișierul \inlinecode{fire\_detection\_system.py} și operează într-un fir de execuție dedicat (\textit{Thread~4}).

Provocarea principală a acestui modul nu constă în detecția propriu-zisă, ci în controlul strict al rezultatelor fals pozitive: focul și, în special, fumul, reprezintă categorii vizuale cu un grad ridicat de ambiguitate (fiind predispuse la confuzii cu pereți gri, formațiuni noroase, abur, reflexii calde sau lumina apusului). O detecție primară obținută de la un model YOLO, dacă ar fi utilizată direct pentru declanșarea unei alarme, ar genera un volum inacceptabil de avertizări false. Din acest motiv, implementarea suprapune peste modelul de bază un \textit{mecanism de filtrare structurat pe cinci niveluri}, care transformă detecțiile candidate în alerte finale exclusiv după ce acestea îndeplinesc, în ordine succesivă, condițiile tuturor filtrelor. Subcapitolele următoare descriu modelul de bază, cele cinci etape de filtrare și, în final, semnificația indicatorilor de stare \inlinecode{confirmed} și \inlinecode{alert} asociați fiecărei detecții.

% [DIAGRAMĂ SUGERATĂ] O diagramă în „cascadă'' (funnel) ar ilustra clar pipeline-ul:
% detecții brute YOLOv10 → [N1] prag de încredere per clasă → [N2] filtru de
% dimensiune → [N3] verificare cromatică HSV → [N4] confirmare temporală IoU
% (≥6 cadre) → [N5] cooldown de alertă pe locație → alertă. La fiecare nivel ar
% apărea ce detecții sunt eliminate.

\subsection{Modelul YOLOv10 pre-antrenat}
\label{subsec:yolov10_preantrenat}

Modelul de bază utilizat este \inlinecode{TommyNgx/YOLOv10-Fire-and-Smoke-Detection}, o variantă a arhitecturii YOLOv10 publicată pe platforma HuggingFace și antrenată specific pentru cele două clase de interes: \inlinecode{Fire} și \inlinecode{Smoke}. Modelul este integrat în sistem în forma sa originală, nefiind necesară o etapă suplimentară de reantrenare (\emph{fine-tuning}), și este instanțiat prin intermediul clasei \inlinecode{YOLO} (din pachetul \inlinecode{ultralytics}) pe baza fișierului de ponderi \inlinecode{models/fire\_and\_smoke/best.pt}. Imediat după încărcare, denumirile claselor sunt normalizate prin convertirea literelor în minuscule (\inlinecode{fire}, \inlinecode{smoke}), asigurând astfel o utilizare consecventă în restul fluxului de prelucrare.

La nivelul fiecărui cadru, metoda \inlinecode{process\_frame()} execută funcția \linebreak \inlinecode{model.predict()} utilizând un prag de confidență egal cu \textit{cea mai mică} valoare dintre pragurile definite per clasă, alături de un prag IoU de \inlinecode{0.45} pentru algoritmul de suprimare a non-maximelor (NMS). Această abordare a fost aleasă în mod deliberat: modelului i se permite să returneze toate detecțiile candidate care depășesc pragul minim de confidență, urmând ca filtrarea riguroasă la nivel de clasă să fie aplicată ulterior, în mod explicit, de către mecanismul de procesare. Această arhitectură oferă un control deplin asupra procesului decizional de declanșare a alertelor.

\subsection{Pipeline-ul de filtrare în cinci niveluri}
\label{subsec:pipeline_filtrare}

Fiecare detecție candidată returnată de YOLOv10 parcurge, în ordine, cinci niveluri de
filtrare. O detecție trebuie să treacă de toate cele cinci niveluri pentru a deveni o alertă;
eșecul la oricare nivel o elimină (sau, în cazul confirmării temporale, o reține ca neconfirmată
încă). Statisticile interne (\inlinecode{rejected\_by\_size}, \inlinecode{rejected\_by\_color}
etc.) contorizează câte detecții sunt respinse la fiecare nivel.

\subsubsection{Praguri de încredere per clasă}
\label{subsubsec:praguri_incredere}

Primul nivel aplică un prag de încredere \textit{diferențiat pe clasă}: \inlinecode{0.40}
pentru \inlinecode{fire} și \inlinecode{0.55} pentru \inlinecode{smoke}. Pragul mai sever pentru
fum reflectă faptul că acesta este semnificativ mai predispus la false pozitive (suprafețe gri,
nori, abur) decât focul, care are o semnătură vizuală mai distinctivă. O detecție a cărei
încredere este sub pragul clasei sale este eliminată imediat.

\subsubsection{Filtre de dimensiune}
\label{subsubsec:filtre_dimensiune}

Al doilea nivel verifică plauzibilitatea dimensiunii bounding box-ului, raportată la aria
cadrului. Sunt respinse atât detecțiile neglijabil de mici (sub \inlinecode{0.3\%} din cadru --
tipic zgomot sau reflexii punctuale), cât și cele excesiv de mari (peste \inlinecode{85\%} din
cadru -- tipic o clasificare greșită care ``acoperă'' aproape întreaga imagine). Doar
detecțiile cu o arie relativă în intervalul $[0{,}003;\ 0{,}85]$ trec mai departe.

\subsubsection{Verificarea cromatică HSV}
\label{subsubsec:verificare_hsv}

Al treilea nivel este o verificare cromatică euristică, realizată în spațiul de culoare HSV de
metoda \inlinecode{\_color\_plausible()}, care exploatează faptul că focul și fumul au semnături
de culoare caracteristice:

\begin{itemize}
    \item \textbf{Foc} -- regiunea trebuie să conțină o fracție semnificativă de pixeli ``calzi'',
    saturați și luminoși: nuanțe (\textit{hue}) în jurul roșu-portocaliu-galben (care se
    încadrează la $h \leq 25$ sau $h \geq 160$, deoarece nuanțele calde se ``înfășoară'' în jurul
    valorii 0), cu saturație $s \geq 90$ și luminozitate $v \geq 120$. Se cere ca cel puțin
    \inlinecode{12\%} din pixeli să fie calzi;
    \item \textbf{Fum} -- regiunea trebuie să fie de saturație redusă și luminozitate moderată:
    cel puțin \inlinecode{55\%} din pixeli trebuie să fie cenușii (saturație $s \leq 60$ și
    luminozitate $v \geq 40$, pentru a respinge zonele complet negre). În plus, deviația
    standard a luminozității trebuie să fie cel puțin $6$, condiție care respinge suprafețele
    perfect uniforme (cer senin, perete), unde fumul real ar prezenta variație.
\end{itemize}

Detecțiile care nu îndeplinesc condiția cromatică a clasei lor sunt respinse, eliminând o mare
parte din falsele pozitive care ar trece de primele două niveluri.

\subsubsection{Confirmarea temporală prin tracking IoU}
\label{subsubsec:confirmare_temporala}

Al patrulea nivel cere persistență în timp: o detecție devine \textit{confirmată} numai dacă
apare în mod consistent pe cel puțin \inlinecode{6} cadre consecutive. Pentru a urmări o aceeași
detecție de la un cadru la altul, metoda \inlinecode{\_update\_frame\_history()} asociază
bounding box-ul curent cu intrările din istoricul de cadre pe baza suprapunerii spațiale
(IoU $\geq 0{,}3$, calculat de \inlinecode{\_compute\_iou()}), restrânsă la aceeași clasă. Dacă
se găsește o potrivire, contorul de cadre consecutive al intrării este incrementat; altfel se
creează o intrare nouă, cu contor 1. Intrările care nu mai sunt observate într-un cadru sunt
expirate (eliminate), întrerupând astfel seria. Acest mecanism elimină detecțiile tranzitorii
de un singur cadru (sclipiri, artefacte de compresie), care nu pot reprezenta un incendiu real.

\subsubsection{Cooldown de alertă pe locație}
\label{subsubsec:cooldown}

Al cincilea nivel previne inundarea cu alerte repetate pentru același eveniment. Odată ce o
detecție este confirmată, se generează o alertă numai dacă a trecut suficient timp de la ultima
alertă pentru \textit{aceeași locație}. Locația este aproximată printr-o cheie grosieră,
construită din clasa detecției și din coordonatele colțului împărțite la 40
(\inlinecode{class\_name\_x1//40\_y1//40}), astfel încât detecții ușor deplasate să fie tratate
ca aceeași sursă. Intervalul minim de cooldown între două alerte la aceeași locație este de
\inlinecode{3} secunde. Astfel, un incendiu persistent generează o alertă inițială și apoi
reamintiri rare, în loc de o avalanșă de alerte la fiecare cadru.

\subsection{Flag-urile \textit{confirmed} și \textit{alert}}
\label{subsec:flaguri}

Fiecare detecție returnată de \inlinecode{process\_frame()} poartă două flag-uri booleene care
codifică rezultatul pipeline-ului de filtrare și care au roluri distincte:

\begin{itemize}
    \item \inlinecode{confirmed} -- adevărat când detecția a trecut de primele patru niveluri și
    a atins pragul de persistență temporală (cel puțin 6 cadre consecutive). O detecție confirmată
    este considerată un eveniment real și este \textit{desenată} pe cadru; candidații neconfirmați
    nu sunt afișați, evitând clipirea vizuală a casetelor nesigure;
    \item \inlinecode{alert} -- adevărat doar în cadrul în care o detecție confirmată trece și de
    nivelul de cooldown, marcând momentul precis în care trebuie generată o \textit{alarmă} (de
    exemplu, înregistrarea în baza de date și notificarea prin WebSocket). Vizual, o detecție cu
    \inlinecode{alert} este desenată cu un chenar mai gros, iar la nivel de cadru se afișează un
    banner ``FIRE ALERT''.
\end{itemize}

Distincția dintre cele două flag-uri separă astfel \textit{starea} (există foc/fum confirmat în
acest cadru?) de \textit{evenimentul} (trebuie declanșată acum o alarmă?), permițând afișarea
continuă a detecțiilor confirmate fără a genera alarme redundante.

\section{Modulul de detecție a armelor (model fine-tuned propriu)}
\label{sec:modul_arme}

Spre deosebire de modulele anterioare, care folosesc modele pre-antrenate preluate ca atare,
modulul de detecție a armelor se bazează pe un model YOLOv8 \textit{fine-tuned} în cadrul
proiectului, pe un set de date construit special pentru această sarcină. Motivul este lipsa
unui model public suficient de robust pentru detecția generică a armelor în context de
supraveghere. Această secțiune descrie întregul proces, de la deciziile de proiectare
(o singură clasă, alegerea modelului de bază), prin construcția și unificarea setului de date,
până la parametrii de antrenare și analiza rezultatelor. Pipeline-ul de antrenare este organizat
în trei scripturi numerotate, rulate în ordine: \inlinecode{01\_download\_datasets.py},
\inlinecode{02\_relabel\_and\_merge.py} și \inlinecode{03\_train\_yolo.py}.

\subsection{Justificarea unei singure clase \textit{weapon}}
\label{subsec:justificare_clasa_unica}

Modelul este antrenat să recunoască o \textit{singură} clasă generică, \inlinecode{weapon},
care acoperă deopotrivă arme de foc (pistoale, puști), arme albe (cuțite, macete) și obiecte
contondente (bâte de baseball). Decizia de a colapsa toate tipurile de arme într-o singură clasă
este motivată practic: în contextul unui sistem de supraveghere, ceea ce contează este
\textit{prezența} unei arme în cadru, nu categoria ei exactă. Această simplificare reduce
problema de clasificare la una de detecție binară (armă / non-armă), permite modelului să
generalizeze mai bine pe exemple noi și elimină confuziile inutile între subtipuri de arme, care
oricum ar declanșa aceeași alertă.

\subsection{Selectarea modelului de bază}
\label{subsec:model_baza}

Modelul de bază ales pentru fine-tuning este \textbfit{YOLOv8m} (varianta \textit{medium}),
pre-antrenat pe datasetul COCO. Această variantă oferă un echilibru bun între acuratețe și
viteză de inferență: este suficient de expresivă pentru a învăța semnătura vizuală variată a
armelor, dar suficient de ușoară pentru a rula în timp real, alături de celelalte modele, pe
același GPU. Fine-tuning-ul pornește de la greutățile COCO (\inlinecode{pretrained=True}),
beneficiind astfel de reprezentările vizuale generice deja învățate, peste care se adaptează
detectorul la noua clasă.

\subsection{Dataseturile utilizate}
\label{subsec:dataseturi_arme}

Întrucât niciun dataset public unic nu acoperea satisfăcător diversitatea de arme, condiții de
iluminare și perspective dorite, au fost combinate două seturi de date publice, descărcate
automat de scriptul \inlinecode{01\_download\_datasets.py}:

\begin{itemize}
    \item \textbf{Dangerous Items} (Zenodo) -- 8.478 de imagini, cu cinci clase originale, toate
    arme: \inlinecode{machete}, \inlinecode{knife}, \inlinecode{baseball\_bat},
    \inlinecode{rifle} și \inlinecode{gun}. Imaginile acoperă scene diverse, de la arme
    prezentate izolat până la arme ținute în mână;
    \item \textbf{SOHAS Weapon Detection} (Roboflow) -- 5.858 de imagini, cu șase clase
    originale: \inlinecode{pistol}, \inlinecode{knife}, \inlinecode{smartphone},
    \inlinecode{billete} (bancnote), \inlinecode{monedero} (portofel) și \inlinecode{tarjeta}
    (card). Ultimele patru clase sunt obiecte non-armă care pot fi ușor confundate cu arme și
    ajută modelul să învețe să le diferențieze.
\end{itemize}

După combinare, datasetul final conține 14.336 de imagini și 12.592 de bounding box-uri,
oferind o diversitate suficientă pentru antrenarea unui detector robust.

\subsection{Preprocesarea și unificarea}
\label{subsec:preprocesare_unificare}

Cele două dataseturi folosesc scheme de etichetare diferite, astfel încât a fost necesară o
etapă de preprocesare pentru unificarea lor într-o structură YOLO coerentă, realizată de scriptul
\inlinecode{02\_relabel\_and\_merge.py}. Operațiile principale sunt:

\begin{itemize}
    \item \textbf{Relabel} -- toate clasele de arme din ambele dataseturi sunt remapate la clasa
    unică \inlinecode{0} (\inlinecode{weapon}). În cazul Zenodo, toate cele cinci clase sunt
    arme, deci toate sunt păstrate; în cazul Roboflow, doar \inlinecode{pistol} și
    \inlinecode{knife} sunt păstrate, iar clasele non-armă (\inlinecode{smartphone},
    \inlinecode{billete}, \inlinecode{monedero}, \inlinecode{tarjeta}) sunt \textit{eliminate}
    din etichete;
    \item \textbf{Merge} -- cele două dataseturi sunt unite într-o singură structură YOLO cu trei
    partiții (\inlinecode{train}, \inlinecode{valid}, \inlinecode{test}); partiția \inlinecode{val}
    a sursei Zenodo este redenumită \inlinecode{valid} pentru consistență;
    \item \textbf{Negative samples} -- imaginile fără adnotări (în care nu apare nicio armă) sunt
    păstrate ca exemple negative. Ele învață modelul ce \textit{nu} este o armă și reduc rata de
    false pozitive la inferență, mai ales pentru obiectele care seamănă cu armele (cazul
    imaginilor cu telefoane sau portofele din Roboflow, ale căror etichete au fost golite).
\end{itemize}

La final, scriptul generează automat un fișier \inlinecode{data.yaml} cu structura standard YOLO,
declarând o singură clasă (\inlinecode{nc: 1}, \inlinecode{names: \{0: weapon\}}).

\subsection{Strategia de denumire (\textit{zen\_} / \textit{rf\_})}
\label{subsec:strategie_denumire}

O problemă subtilă la unirea a două dataseturi este coliziunea numelor de fișiere: ambele surse
pot conține, de exemplu, o imagine numită \inlinecode{0001.jpg}, care s-ar suprascrie reciproc la
copierea în directorul comun. Pentru a o evita, fiecărei imagini și etichete i se adaugă un prefix
în funcție de sursă: \inlinecode{zen\_} pentru imaginile din Zenodo și \inlinecode{rf\_} pentru
cele din Roboflow. Astfel, perechile imagine--etichetă rămân corelate (același prefix și același
nume de bază), iar fișierele celor două surse coexistă fără pierderi în arborele
\inlinecode{train}/\inlinecode{valid}/\inlinecode{test} comun.

\subsection{Parametrii de antrenare}
\label{subsec:parametri_antrenare_arme}

Antrenarea propriu-zisă, realizată de scriptul \inlinecode{03\_train\_yolo.py}, a folosit
următorii parametri principali:

\begin{itemize}
    \item \textbf{Epoci}: 100 de epoci planificate, cu \textit{early stopping}
    (\inlinecode{patience=20}) -- oprire automată dacă metrica de validare nu se mai îmbunătățește
    timp de 20 de epoci;
    \item \textbf{Batch size}: 16 imagini per batch;
    \item \textbf{Rezoluția de intrare}: $640 \times 640$ pixeli, valoarea standard pentru YOLOv8;
    \item \textbf{Optimizator}: AdamW, cu learning rate inițial $lr_0 = 0{,}001$ și factor final
    $lr_f = 0{,}01$, conform unei programări cosinus a ratei de învățare
    (\inlinecode{cos\_lr=True}), și \inlinecode{weight\_decay} $= 0{,}0005$;
    \item \textbf{Warmup}: 3 epoci de warmup pentru stabilizarea gradienților la începutul
    antrenării.
\end{itemize}

Scriptul include și o rutină de selecție automată a GPU-ului cu cea mai multă memorie liberă
(\inlinecode{pick\_gpu()}, via \inlinecode{pynvml} sau, în lipsa lui, \inlinecode{nvidia-smi}).
Antrenarea a fost realizată pe un GPU NVIDIA A100.

\subsection{Augmentările folosite}
\label{subsec:augmentari}

Pentru a crește robustețea modelului și a reduce overfitting-ul, în timpul antrenării au fost
aplicate multiple transformări de augmentare a datelor:

\begin{itemize}
    \item \textbf{Variații cromatice HSV} -- nuanță ($\pm 0{,}015$), saturație ($\pm 0{,}7$) și
    luminozitate ($\pm 0{,}4$), pentru a simula condiții de iluminare diferite;
    \item \textbf{Transformări geometrice} -- rotație până la $10^{\circ}$, translație de
    $10\%$, scalare de $50\%$ și oglindire orizontală (\inlinecode{fliplr=0.5});
    \item \textbf{Mosaic} (\inlinecode{mosaic=1.0}) -- combinarea a patru imagini într-una
    singură, expunând modelul la contexte și scale variate;
    \item \textbf{MixUp} (\inlinecode{mixup=0.1}) -- amestecarea a două imagini, ca regularizare
    suplimentară.
\end{itemize}

\subsection{Rezultatele antrenării}
\label{subsec:rezultate_arme}

Antrenarea a parcurs toate cele 100 de epoci (early stopping nu s-a declanșat). Metricile pe
partiția de validare la finalul antrenării sunt sintetizate în tabelul~\ref{tab:rezultate_arme}.

\begin{table}[h]
    \centering
    \caption{Metricile finale de validare ale modelului de detecție a armelor.}
    \label{tab:rezultate_arme}
    \begin{tabular}{l c}
        \hline
        \textbf{Metrică} & \textbf{Valoare} \\
        \hline
        Precizie (\textit{precision}) & $0{,}937$ \\
        Sensibilitate (\textit{recall}) & $0{,}894$ \\
        mAP@$0{,}5$ & $0{,}947$ \\
        mAP@$0{,}5{:}0{,}95$ & $0{,}700$ \\
        \hline
    \end{tabular}
\end{table}

Modelul atinge o precizie de $93{,}7\%$ și o sensibilitate de $89{,}4\%$, cu un mAP@$0{,}5$ de
$0{,}947$. Valoarea mAP@$0{,}5{:}0{,}95$ (media mAP peste praguri IoU de la $0{,}5$ la $0{,}95$)
este de $0{,}700$, indicând o localizare bună, deși mai puțin precisă la pragurile IoU foarte
stricte -- comportament tipic pentru detecția unor obiecte cu forme și orientări foarte variate,
cum sunt armele. Pe lângă greutăți, antrenarea a produs și artefacte de diagnostic (curbele
precizie-sensibilitate, matricea de confuzie, curbele de pierdere), salvate în directorul de
ieșire al rulării.

\subsection{Discuție asupra rezultatelor}
\label{subsec:discutie_arme}

Valorile obținute sunt solide pentru o sarcină dificilă: detecția unei clase vizual eterogene
(o armă de foc, un cuțit și o bâtă de baseball au foarte puține trăsături comune), pe un dataset
combinat din surse cu stiluri de adnotare diferite. Diferența dintre precizie ($93{,}7\%$) și
sensibilitate ($89{,}4\%$) arată că modelul este ușor mai conservator: ratează unele arme, dar
generează relativ puține detecții false -- un compromis dezirabil pentru a evita alarmele false
într-un sistem de securitate.

Acest comportament este întărit la inferență de o decizie de proiectare la nivelul modulului de
rulare (\inlinecode{WeaponDetectionSystem}): deși modelul produce candidați peste un prag de
încredere permisiv ($0{,}3$), o detecție este reținută doar dacă depășește un post-filtru strict
de încredere ($0{,}7$), iar alerta este confirmată temporal pe mai multe cadre consecutive
(mecanism analog celui de la detecția de foc, secțiunea~\ref{subsec:confirmare_temporala}).
Astfel, sensibilitatea brută a modelului este deliberat schimbată pe o precizie operațională mai
mare, prioritizând reducerea falselor alarme. Aportul exemplelor negative (telefoane, portofele,
carduri din SOHAS, cu etichete golite) se reflectă tocmai în această precizie ridicată, modelul
fiind expus explicit la obiecte care pot fi confundate cu arme.

\section{Modulul de recunoaștere a acțiunilor umane (fine-tuning SlowFast)}
\label{sec:modul_actiuni}

Cel mai complex modul de inteligență artificială al sistemului este cel de recunoaștere a
acțiunilor umane (HAR -- \textit{Human Action Recognition}), care clasifică \textit{secvențe
video} (nu cadre izolate) în trei categorii comportamentale: \inlinecode{normal} (activitate
obișnuită), \inlinecode{fight} (lupte, altercații fizice) și \inlinecode{vandalism} (acte de
distrugere intenționată a proprietății). Spre deosebire de detectoarele de obiecte, care
operează pe un singur cadru, recunoașterea acțiunilor necesită modelarea dimensiunii temporale
-- aceeași imagine poate aparține atât unei îmbrățișări, cât și unei altercații, distincția
fiind dată de mișcare.

Acest modul reprezintă, alături de detectorul de arme, una dintre cele două componente
antrenate în cadrul proiectului. Antrenarea (fine-tuning) a fost realizată în proiectul separat
\inlinecode{security\_system}, iar modelul rezultat este consumat de aplicația principală ca un
fișier checkpoint (\inlinecode{best\_model.pth}), încărcat la rulare de clasa
\inlinecode{HumanActionRecognitionSystem} din \inlinecode{har\_system.py} (\textit{Thread~5}).
Subcapitolele care urmează descriu alegerea arhitecturii, construcția setului de date (inclusiv
crearea de la zero a unui dataset propriu pentru clasa \inlinecode{vandalism}), procesul de
fine-tuning, parametrii și rezultatele, iar la final modul de integrare în aplicație.

\subsection{Justificarea alegerii SlowFast R50}
\label{subsec:justificare_slowfast}

Arhitectura aleasă este \textbfit{SlowFast R50}, dezvoltată de Facebook AI Research, ale cărei
fundamente au fost prezentate în secțiunea~\ref{subsec:slowfast}. Ea se bazează pe două căi de
procesare paralele: o cale \textit{Slow}, care procesează puține cadre la rezoluție temporală
redusă și captează informația semantică spațială, și o cale \textit{Fast}, care procesează multe
cadre și captează dinamica mișcării, cele două fiind unite prin conexiuni laterale. Această
separare explicită a aspectului spațial de cel temporal o face deosebit de potrivită pentru
distingerea acțiunilor pe baza mișcării -- exact problema de față.

Varianta R50 (cu backbone ResNet-50) oferă un compromis bun între capacitate și cost
computațional, permițând inferența în timp real alături de celelalte modele. În plus, modelul este
disponibil pre-antrenat pe datasetul Kinetics-400 (prin PyTorchVideo / PyTorch Hub), ceea ce
oferă un punct de pornire solid: reprezentările video generice învățate pe 400 de clase de acțiuni
pot fi adaptate prin fine-tuning la cele trei clase ale problemei, cu un set de date mult mai mic
decât ar fi necesar pentru antrenarea de la zero.

\subsection{Construcția setului de date}
\label{subsec:constructie_dataset}

Setul de date pentru fine-tuning a fost construit din trei surse, câte una pentru fiecare clasă.
Structura este de tip folder-per-clasă (\inlinecode{data/train/<clasă>/} și
\inlinecode{data/test/<clasă>/}), iar clasele sunt descoperite automat din numele folderelor de
către \inlinecode{configs/config.py}, cu convenția ca eticheta \inlinecode{normal} să primească
întotdeauna indexul 0.

\subsubsection{Clasa \textit{normal}}
\label{subsubsec:clasa_normal}

Clasa \inlinecode{normal} conține videoclipuri de supraveghere cu activitate cotidiană, fără
incidente, preluate din datasetul public \textbf{RWF-2000} (disponibil pe Kaggle). Aceste
secvențe oferă modelului o referință pentru ceea ce înseamnă comportament obișnuit, esențială
pentru a delimita clasele de interes (\inlinecode{fight}, \inlinecode{vandalism}) de fundalul
normal.

\subsubsection{Clasa \textit{fight}}
\label{subsubsec:clasa_fight}

Clasa \inlinecode{fight} conține videoclipuri cu lupte și altercații fizice, preluate tot din
datasetul \textbf{RWF-2000}, care include scene de confruntare fizică între persoane, capturate de
camere de supraveghere. Folosirea aceluiași dataset pentru clasele \inlinecode{normal} și
\inlinecode{fight} asigură un stil vizual consistent (rezoluție, unghiuri, calitate a imaginii)
între cele două, astfel încât modelul să învețe diferența de \textit{acțiune}, nu de domeniu vizual.

\subsubsection{Clasa \textit{vandalism} -- dataset propriu}
\label{subsubsec:clasa_vandalism}

\paragraph{Justificarea creării unui dataset propriu.}
\label{par:justificare_dataset}
Spre deosebire de primele două clase, pentru \inlinecode{vandalism} nu există un dataset public
adecvat de videoclipuri de supraveghere. Din acest motiv, a fost necesară construirea unui set de
date propriu, de la zero, printr-un proces în trei etape descris în continuare.

\paragraph{Etapa 1: Scraping automat YouTube.}
\label{par:scraping_youtube}
A fost dezvoltat un script Python dedicat care folosește biblioteca \inlinecode{yt-dlp} pentru a
căuta și colecta automat link-uri către videoclipuri candidate de pe YouTube. Scriptul generează
sute de interogări de căutare variate, în mai multe limbi, acoperind acte specifice de vandalism
(graffiti, spargere de geamuri, distrugere de mașini, intrare prin efracție etc.) combinate cu
termeni asociați camerelor de supraveghere. În urma acestui proces au fost colectate aproximativ
4.800 de link-uri către potențiale videoclipuri.

\paragraph{Etapa 2: Filtrare manuală.}
\label{par:filtrare_manuala}
Toate cele aproximativ 4.800 de videoclipuri au fost vizualizate și evaluate manual, pentru a
selecta doar clipurile care conțin efectiv acte de vandalism filmate de camere de supraveghere,
eliminându-le pe cele irelevante. În urma filtrării au rămas peste 600 de videoclipuri valide.
Această etapă, deși laborioasă, este esențială pentru calitatea datasetului: garantează că modelul
învață din exemple reale și relevante.

\paragraph{Etapa 3: Editare în OpenShot.}
\label{par:editare_openshot}
Videoclipurile selectate au fost editate individual cu software-ul OpenShot Video Editor, pentru
a extrage doar segmentele relevante. Fiecare clip a fost decupat astfel încât să conțină secvența
propriu-zisă a actului de vandalism, eliminând porțiunile irelevante (intro-uri, text suprapus,
secvențe fără activitate).

\paragraph{Setul final.}
\label{par:set_final}
Setul de date final pentru clasa \inlinecode{vandalism} conține \textbf{614 videoclipuri} cu acte
de vandalism, acoperind o gamă largă de acțiuni: graffiti / street art, spargere de geamuri și
vitrine, distrugere de autoturisme (zgâriere, lovire), intrare prin efracție, distrugere de
mobilier urban, aruncare cu obiecte și alte forme de distrugere intenționată a proprietății.

\subsection{Pregătirea modelului pentru fine-tuning}
\label{subsec:pregatire_finetuning}

Pregătirea modelului pornește de la SlowFast R50 pre-antrenat pe Kinetics-400, al cărui cap de
clasificare original -- un strat \inlinecode{Linear(2304, 400)} corespunzător celor 400 de clase
Kinetics -- este înlocuit cu un cap nou, adaptat problemei: o secvență
\inlinecode{Dropout(0.5)} urmată de un \inlinecode{Linear(2304, 3)}, pentru cele trei clase. Această
înlocuire este realizată în \inlinecode{build\_slowfast\_model()} din
\inlinecode{models/slowfast\_model.py}.

\subsubsection{Formatul intrării (contractul SlowFast)}
Fiecare videoclip este transformat într-o pereche de tensori, câte unul pentru fiecare cale:
calea Slow primește \inlinecode{8} cadre, iar calea Fast \inlinecode{32} de cadre, ambele decupate
spațial la $224 \times 224$ pixeli. Cadrele sunt eșantionate uniform dintr-un clip
corespunzător unei durate de \inlinecode{2.0} secunde. Aceste valori sunt fixate în
\inlinecode{configs/config.py} și trebuie respectate identic la inferență (în aplicație), altfel
modelul ar primi intrări incompatibile cu cele pe care a fost antrenat.

\subsection{Parametrii de antrenare}
\label{subsec:parametri_actiuni}

Antrenarea, implementată în \inlinecode{train.py}, folosește o strategie de fine-tuning în
\textit{două faze}, gândită pentru a adapta modelul la noile clase fără a distruge reprezentările
pre-antrenate:

\begin{itemize}
    \item \textbf{Faza 1 -- backbone înghețat} (primele 5 epoci): toți parametrii backbone-ului
    sunt înghețați (\inlinecode{requires\_grad=False}), antrenându-se \textit{doar} capul de
    clasificare nou. Astfel, capul aleator inițializat se stabilizează fără a perturba
    caracteristicile vizuale învățate pe Kinetics-400;
    \item \textbf{Faza 2 -- fine-tuning complet}: backbone-ul este dezghețat și întregul model este
    antrenat, dar cu \textit{rate de învățare diferențiate} -- backbone-ul primește o rată de 10
    ori mai mică decât capul (\inlinecode{lr\_backbone\_factor=0.1}), pentru a-l ajusta fin fără a-l
    deteriora.
\end{itemize}

Alți parametri și tehnici folosite, conform \inlinecode{config.py} și \inlinecode{train.py}:

\begin{itemize}
    \item \textbf{Optimizator}: AdamW, rată de învățare $10^{-3}$, \inlinecode{weight\_decay}
    $= 10^{-4}$, cu o programare de tip \inlinecode{StepLR} (reducere de 10$\times$ la fiecare 10
    epoci);
    \item \textbf{Funcția de pierdere}: \inlinecode{CrossEntropyLoss} cu \textit{label smoothing}
    de $0{,}1$, pentru a reduce supraîncrederea modelului;
    \item \textbf{Dezechilibrul de clase}: tratat printr-un \inlinecode{WeightedRandomSampler}, care
    eșantionează clasele invers proporțional cu frecvența lor, astfel încât fiecare clasă să fie
    reprezentată echilibrat în batch-uri;
    \item \textbf{Tehnici de stabilizare}: antrenare în precizie mixtă (AMP, prin
    \inlinecode{GradScaler} și \inlinecode{autocast}), \textit{gradient clipping} la norma maximă
    $5{,}0$ și \textit{early stopping} (răbdare de 7 epoci) bazat pe macro-F1 de validare;
    \item \textbf{Augmentare spațială}: oglindire orizontală, \textit{color jitter} și
    \textit{random crop}.
\end{itemize}

Din cele 30 de epoci planificate, antrenarea s-a oprit la epoca 28 prin early stopping, cele mai
bune rezultate (cele salvate ca \inlinecode{best\_model.pth}) fiind obținute la epoca 21.
Antrenarea a folosit 262 de batch-uri pe epocă la antrenare și 56 la validare, cu o durată medie
de aproximativ 15 minute per epocă.

\subsection{Rezultatele}
\label{subsec:rezultate_actiuni}

Metricile globale obținute pe setul de validare de către modelul cel mai bun (epoca 21) sunt
sintetizate în tabelul~\ref{tab:rezultate_actiuni}. Pierderea finală de antrenare a fost de
$0{,}3065$, iar cea de validare de $0{,}4306$.

\begin{table}[h]
    \centering
    \caption{Metricile globale de validare ale modelului de recunoaștere a acțiunilor.}
    \label{tab:rezultate_actiuni}
    \begin{tabular}{l c}
        \hline
        \textbf{Metrică} & \textbf{Valoare} \\
        \hline
        Acuratețe (\textit{accuracy}) & $0{,}9433$ \\
        Macro Precision & $0{,}9478$ \\
        Macro Recall & $0{,}9436$ \\
        Macro F1-Score & $0{,}9456$ \\
        \hline
    \end{tabular}
\end{table}

Modelul atinge o acuratețe globală de $94{,}33\%$ și un macro F1-score de $94{,}56\%$, valori
foarte bune pentru o problemă de clasificare a acțiunilor pe trei clase, în condiții de
supraveghere reale.

\subsection{Analiza per clasă}
\label{subsec:analiza_per_clasa}

Rezultatele detaliate pe fiecare clasă sunt prezentate în tabelul~\ref{tab:per_clasa_actiuni}.

\begin{table}[h]
    \centering
    \caption{Metricile per clasă ale modelului de recunoaștere a acțiunilor.}
    \label{tab:per_clasa_actiuni}
    \begin{tabular}{l c c c c}
        \hline
        \textbf{Clasă} & \textbf{Precision} & \textbf{Recall} & \textbf{F1-Score} & \textbf{Suport} \\
        \hline
        \inlinecode{normal}    & $0{,}920$ & $0{,}947$ & $0{,}933$ & $169$ \\
        \inlinecode{fight}     & $0{,}942$ & $0{,}935$ & $0{,}939$ & $155$ \\
        \inlinecode{vandalism} & $0{,}982$ & $0{,}949$ & $0{,}965$ & $117$ \\
        \hline
    \end{tabular}
\end{table}

Se remarcă faptul că tocmai clasa \inlinecode{vandalism} -- cea pentru care a fost construit
datasetul propriu -- obține cele mai bune rezultate, cu o precizie de $98{,}2\%$ și un F1-score de
$0{,}965$. Acest rezultat validează efortul de creare a datasetului custom: clipurile selectate
și editate manual sunt suficient de curate și de reprezentative pentru ca modelul să învețe
robust această clasă. Clasele \inlinecode{normal} și \inlinecode{fight} obțin de asemenea scoruri
ridicate (F1 de $0{,}933$, respectiv $0{,}939$), confuziile reziduale fiind, în mod așteptat,
preponderent între ele -- ambele provin din același dataset RWF-2000 și pot conține secvențe
ambigue la granița dintre activitate agitată și altercație.

\subsubsection{Integrarea în aplicație}
La rulare, modelul antrenat este folosit de \inlinecode{HumanActionRecognitionSystem}, care nu
poate clasifica un singur cadru, ci are nevoie de o secvență. De aceea, sistemul acumulează
cadrele primite într-un \textit{ring buffer} (\inlinecode{deque}) și rulează inferența doar o dată
la fiecare \inlinecode{30} de cadre (\inlinecode{clip\_interval\_frames}), atunci când bufferul
conține destule cadre. La fiecare inferență, se eșantionează uniform până la 64 de cadre din
buffer, se construiește perechea slow/fast (8/32 cadre, $224 \times 224$) exact ca la antrenare, se
aplică \inlinecode{softmax} peste logits și se reține clasa cu probabilitatea maximă. O acțiune
non-normală este raportată ca alertă doar dacă încrederea depășește pragul
\inlinecode{confidence\_threshold} ($0{,}5$), cu un \textit{cooldown} de alertă de 5 secunde pentru
a evita notificările repetate. Predicția curentă este păstrată între inferențe, astfel încât
adnotarea afișată pe flux rămâne coerentă chiar și pe cadrele dintre două rulări de inferență.

\section{Fuziunea rezultatelor pe cadru}
\label{sec:fuziune}

Cele cinci module de inteligență artificială descrise anterior nu rulează izolat, ci sunt
orchestrate de aplicația principală (\inlinecode{app.py}) într-un \textit{pipeline
multi-threaded} care procesează același flux video în paralel și apoi recombină rezultatele
într-un singur cadru adnotat. Această secțiune descrie modul în care detecțiile provenite de la
modele diferite, care rulează pe fire de execuție separate și se termină la momente diferite,
sunt sincronizate, suprapuse pe cadru, transformate în alarme și difuzate către clienți.

Decizia arhitecturală fundamentală este aceea de a folosi \textit{fire de execuție}
(\inlinecode{threading}) cu cozi (\inlinecode{queue.Queue}), nu programare asincronă: apelurile
de inferență ale modelelor sunt operații blocante, intensive computațional, care nu trebuie să
blocheze bucla de evenimente a serverului FastAPI. Pipeline-ul folosește în total
\textbf{șapte fire de execuție}: captură (Thread~1), recunoaștere facială cu demografie
(Thread~2), plăcuțe (Thread~3), foc și fum (Thread~4), acțiuni umane (Thread~5), arme (Thread~6)
și un fir final de fuziune (Thread~7).

% [DIAGRAMĂ SUGERATĂ] O diagramă a arhitecturii multi-threaded ar fi foarte utilă aici:
% Thread 1 (captură) → dispecerizare în 5 cozi paralele → Thread 2–6 (cele cinci
% module AI, fiecare cu coada lui de intrare și de rezultate) → Thread 7 (fuziune,
% sincronizare după frame_id) → encodare JPEG → broadcast WebSocket către clienți.

\subsection{Dispecerizarea cadrelor și procesarea în paralel}
\label{subsec:dispecerizare}

Firul de captură (Thread~1) citește cadre de la toate camerele active (camera laptop-ului și
eventualele camere IP). Fiecărui cadru i se atribuie un identificator unic și monoton crescător,
\inlinecode{frame\_id}, care va servi drept cheie de sincronizare pe tot parcursul pipeline-ului.
Cadrul este apoi \textit{dispecerizat} de metoda \inlinecode{\_dispatch\_frame()}, care îi pune
câte o copie în coada de intrare a fiecărui modul activat (\inlinecode{face\_processing\_queue},
\inlinecode{plate\_processing\_queue} etc.) și, separat, în coada de cadre brute
(\inlinecode{raw\_frame\_queue}), care păstrează imaginea originală pentru recompunere.

Fiecare modul AI rulează în propriul fir, consumând cadre din coada sa, executând inferența și
publicând rezultatele în coada sa de ieșire (\inlinecode{face\_results\_queue} etc.). Cum cozile
au capacitate mărginită (\inlinecode{maxsize=10}), atunci când un modul lent nu mai poate prelua
cadre, dispecerul nu blochează, ci sare peste cadru (incrementând un contor de
\inlinecode{queue\_skips} / \inlinecode{frames\_dropped}). Astfel, un model mai lent degradează
elegant rata de cadre, fără a opri întregul sistem.

\subsection{Sincronizarea rezultatelor după \textit{frame\_id}}
\label{subsec:sincronizare}

Firul de fuziune (Thread~7) are sarcina dificilă de a recombina rezultate care sosesc
\textit{asincron}: modulul de plăcuțe poate termina cadrul $N$ înainte ca modulul facial să
termine cadrul $N-2$. Pentru aceasta, firul menține un dicționar \inlinecode{pending\_results},
cu câte o intrare pentru fiecare tip de rezultat (\inlinecode{frames}, \inlinecode{face},
\inlinecode{plate}, \inlinecode{fire}, \inlinecode{har}, \inlinecode{weapon}), indexată după
\inlinecode{frame\_id}. La fiecare iterație, firul golește toate cozile de rezultate disponibile
în acest dicționar, apoi încearcă să \textit{asambleze} fiecare cadru pentru care a sosit imaginea
brută.

Un cadru este considerat gata de asamblare atunci când rezultatele modulelor \textit{esențiale}
activate (recunoaștere facială și plăcuțe) au sosit. Rezultatele modulelor de foc, acțiuni și arme
sunt tratate ca \textit{opționale}: ele sunt suprapuse dacă au sosit, dar absența lor nu blochează
asamblarea. Pentru a evita blocajul atunci când un rezultat nu mai sosește (de exemplu, din cauza
unui cadru sărit), există un mecanism de \textit{timeout}: dacă un cadru devine prea vechi
(diferența dintre \inlinecode{frame\_id}-ul curent și al lui depășește 60), el este asamblat cu
rezultatele disponibile la acel moment, fără a mai aștepta. În plus, o rutină de curățare
(\inlinecode{\_cleanup\_old\_results()}) elimină periodic intrările mai vechi de 90 de cadre,
pentru ca dicționarul să nu crească nelimitat în caz de desincronizare.

\subsection{Compunerea cadrului adnotat}
\label{subsec:compunere}

Odată ce un cadru este gata de asamblare, adnotările celor cinci module sunt desenate
\textit{în straturi} peste imagine, într-o ordine fixă. Punctul de plecare este cadrul deja
adnotat de modulul facial (care conține chenarele fețelor și etichetele demografice), peste care
se aplică succesiv: chenarele plăcuțelor (\inlinecode{plate\_system.draw\_outputs()}), detecțiile
de foc confirmate (\inlinecode{fire\_system.draw\_detections()}), eticheta acțiunii recunoscute
(\inlinecode{har\_system.draw\_detections()}) și, la final, chenarele armelor
(\inlinecode{weapon\_system.draw\_detections()}). Rezultatul este un singur cadru
(\inlinecode{final\_frame}) care concentrează vizual toate detecțiile celor cinci module pentru
acel moment de timp.

\subsection{Generarea alarmelor și deduplicarea}
\label{subsec:alarme}

În aceeași etapă de fuziune, detecțiile relevante sunt transformate în două tipuri de înregistrări
persistente: \textit{jurnale de detecție} (un istoric complet, scris prin
\inlinecode{\_log\_detection\_if\_needed()}) și \textit{alarme} propriu-zise (evenimente care
necesită atenția operatorului). Sunt considerate demne de alarmă: persoanele necunoscute, focul
confirmat cu \inlinecode{alert}, acțiunile non-\inlinecode{normal}, armele detectate, vehiculele
neautorizate (plăcuțe lizibile care nu sunt pe lista albă) și prezența unei persoane neautorizate
într-o zonă restricționată.

Problema centrală a generării alarmelor este evitarea inundării operatorului cu duplicate: un foc
care arde 10 secunde nu trebuie să genereze o alarmă la fiecare cadru. Mecanismul de deduplicare,
implementat în \inlinecode{\_create\_alarm\_if\_needed()}, folosește o cheie compusă din cameră,
tip și \textit{subiect} (\inlinecode{camera\_id:alarm\_type:dedup\_subject}) și un cooldown de 30
de secunde per cheie. Componenta de subiect este esențială: ea permite ca alarme distincte de
\textit{același} tip să coexiste -- de exemplu, două persoane necunoscute diferite (diferențiate
prin poziție), un cuțit și un pistol, sau două plăcuțe neautorizate diferite -- fără a se
contopi într-o singură alarmă. Pe lângă bariera din memorie, există și o verificare de rezervă în
baza de date (\inlinecode{get\_recent\_alarm()}), care previne avalanșa de alarme duplicate
imediat după o repornire a serverului, când starea de cooldown din memorie este goală.

\subsection{Difuzarea către clienți prin WebSocket}
\label{subsec:difuzare}

Ultimul pas al fuziunii este livrarea cadrului asamblat către interfețele clienților. Cadrul
\inlinecode{final\_frame} este redimensionat (la maximum 800~px lățime), codat JPEG cu o calitate
de 75 și convertit în Base64 prin \inlinecode{\_encode\_and\_queue\_frame()}. Acesta este împachetat
într-un mesaj JSON care conține, pe lângă imagine, listele de rezultate serializate ale tuturor
modulelor (\inlinecode{face\_results}, \inlinecode{plate\_results}, \inlinecode{fire\_results},
\inlinecode{har\_results}, \inlinecode{weapon\_results}), identificatorul camerei, un
\textit{timestamp} și starea (activat/dezactivat) a fiecărui modul.

Difuzarea propriu-zisă (\inlinecode{\_broadcast\_frame()}) folosește un model de tip
\textit{fan-out} cu cozi independente per client: fiecare client conectat prin WebSocket are
propria coadă de ieșire, iar firul de fuziune publică fiecare cadru procesat în toate cozile.
Această izolare este importantă -- un client lent își aruncă doar propriul cadru cel mai vechi
atunci când coada lui se umple, fără a ``fura'' cadre de la ceilalți clienți sau a-i încetini.
Astfel, fluxul video adnotat ajunge la fiecare operator în mod independent, încheind ciclul de
procesare care a pornit de la captura cadrului brut.

\end{document}
